\section{Discussion}
\label{sec:discussion}

\subsection{Hypothesis Evaluation}

The primary hypothesis is \textbf{confirmed}. Recursive proof
composition reduces per-enterprise verification gas by approximately
$N$-fold or better for all three aggregation strategies tested.

For the primary target of $N{=}8$ enterprises:

\begin{itemize}
  \item \textbf{Folding + Groth16 decider:} $15.3\times$ gas reduction
    (3.36M $\to$ 220K total; 420K $\to$ 27.5K per enterprise). This
    exceeds the predicted $N$-fold because the Groth16 decider at 220K
    gas is cheaper than the original halo2-KZG verification at 420K
    gas, yielding approximately $2N$-fold savings.
  \item \textbf{Binary tree accumulation:} $8.0\times$ gas reduction
    (3.36M $\to$ 420K), exactly $N$-fold as the final proof uses the
    same halo2-KZG verification.
  \item \textbf{SnarkPack:} $7.5\times$ gas reduction (3.36M $\to$
    450K), slightly below $N$-fold due to $O(\log N)$ overhead in the
    batch proof size.
\end{itemize}

The null hypothesis (linear or worse scaling of overhead) is
\textbf{rejected}. All three strategies demonstrate sub-linear
aggregation overhead: folding time grows as $O(N)$ with small constant
(${\sim}250$\,ms/fold), binary tree time grows as $O(N \log N)$, and
SnarkPack time grows as $O(N \log N)$ with very small constant.

\subsection{Success Criteria Assessment}

All four pre-registered success criteria are met:

\begin{enumerate}
  \item Gas reduction $\geq 4\times$ at $N{=}8$: \textbf{Achieved}
    ($15.3\times$ for folding, $8.0\times$ for binary tree).
  \item Aggregation time $< 30$\,s at $N{=}8$: \textbf{Achieved}
    (11.75\,s for folding, 200\,ms for SnarkPack).
  \item Proof size $< 2$\,KB: \textbf{Achieved} (128\,B for folding,
    640\,B for binary tree, 592\,B for SnarkPack).
  \item Soundness preservation: \textbf{Achieved} by construction
    (recursive soundness of SNARKs, formalized in
    Section~\ref{sec:invariants}).
\end{enumerate}

\subsection{Aggregation Invariants}
\label{sec:invariants}

The experiments reveal four invariants that must hold in any production
aggregation deployment:

\begin{description}
  \item[INV-AGG-1 (Aggregation Soundness).] An aggregated proof is valid
    if and only if \emph{all} $N$ component proofs are valid. No subset
    of valid proofs can make an aggregated proof valid if any component
    is invalid.

  \item[INV-AGG-2 (Independence Preservation).] Enterprise provers do
    not need to coordinate. Failure of one enterprise's proof does not
    invalidate other enterprises' proofs. The aggregator handles partial
    sets gracefully by aggregating whatever valid proofs are available.

  \item[INV-AGG-3 (Order Independence).] The aggregated proof is
    deterministic regardless of the order in which enterprise proofs are
    presented to the aggregator. This is required for reproducibility
    and auditing.

  \item[INV-AGG-4 (Gas Monotonicity).] The amortized per-enterprise gas
    cost strictly decreases as $N$ increases. Adding an enterprise to
    the aggregation set never increases per-enterprise cost.
\end{description}

These invariants will be formalized in TLA+ by The Logicist (RU-L10
downstream, checklist item~38) and verified in Coq by The Prover.

\subsection{Strategy Comparison}

Table~\ref{tab:strategy-comparison} presents a multi-criteria comparison
of the three strategies.

\begin{table}[htbp]
\centering
\caption{Strategy Comparison (Multi-Criteria)}
\label{tab:strategy-comparison}
\begin{tabular}{@{}p{2.2cm}p{1.5cm}p{1.5cm}p{1.7cm}@{}}
\toprule
\textbf{Criterion} & \textbf{Binary Tree} & \textbf{Folding+G16} & \textbf{SnarkPack} \\
\midrule
Gas ($N{=}8$)           & 420K         & \textbf{220K} & 450K \\
Gas/Enterprise          & 52.5K        & \textbf{27.5K}& 56.3K \\
Savings ($N{=}8$)       & 8.0$\times$  & \textbf{15.3$\times$} & 7.5$\times$ \\
Agg.\ time ($N{=}8$)    & 180\,s       & 11.75\,s      & \textbf{0.2\,s} \\
Final proof size        & 640\,B       & \textbf{128\,B}& 592\,B \\
Memory scaling          & $O(N)$       & \textbf{$O(1)$}& $O(N)$ \\
Trusted setup           & None extra   & Per-circuit    & None extra \\
Production maturity     & \textbf{High}& Low           & Medium \\
halo2 compatibility     & \textbf{Direct}& Good        & Groth16 only \\
\bottomrule
\end{tabular}
\end{table}

\textbf{Folding + Groth16 decider} dominates on gas efficiency, proof
size, memory scaling, and aggregation time. Its weaknesses are the
requirement for a per-circuit Groth16 trusted setup for the decider and
the immaturity of the Sonobe library (experimental, not audited).

\textbf{Binary tree accumulation} dominates on production maturity and
tooling (Scroll processes millions of dollars daily using this
approach). Its weakness is aggregation time (180\,s at $N{=}8$), though
this is not user-facing latency---aggregation happens asynchronously
before L1 submission.

\textbf{SnarkPack} provides the fastest aggregation (200\,ms) but
requires Groth16 inputs (not directly compatible with halo2-KZG
proofs without a wrapper) and offers lower gas savings than folding.

\subsection{Scaling Projections}

Extrapolating from measured data to larger $N$:

\begin{itemize}
  \item \textbf{Folding + Groth16:} At $N{=}32$, predicted gas savings
    of $61.1\times$ (13.44M $\to$ 220K), aggregation time ${\sim}17$\,s,
    memory constant at 1.6\,GB.
  \item \textbf{Binary tree:} At $N{=}32$, predicted savings of
    $32\times$ (13.44M $\to$ 420K), aggregation time ${\sim}300$\,s.
  \item \textbf{SnarkPack:} At $N{=}32$, predicted savings of
    ${\sim}24\times$ (13.44M $\to$ 550K), aggregation time ${\sim}800$\,ms.
\end{itemize}

The constant-memory property of folding schemes (1.6\,GB regardless of
$N$, per Nova~\cite{nova}) makes them the only viable strategy for very
large $N$ ($> 100$ enterprises). Binary tree accumulation memory grows
as $O(N)$, reaching ${\sim}32$\,GB at $N{=}100$.

\subsection{Phased Deployment Strategy}
\label{sec:deployment}

We recommend a two-phase deployment strategy that balances engineering
risk against gas efficiency:

\textbf{Phase 1 (Near-term):} Binary tree accumulation via the halo2
snark-verifier library.

\begin{itemize}
  \item Production-proven at Scroll and Axiom on Ethereum mainnet.
  \item $8\times$ gas savings for $N{=}8$ enterprises.
  \item Higher aggregation time (180\,s) is acceptable as aggregation is
    not on the critical path for user-facing latency.
  \item Uses the same universal SRS as enterprise provers---no additional
    trusted setup.
\end{itemize}

\textbf{Phase 2 (Production):} ProtoGalaxy folding + Groth16 decider.

\begin{itemize}
  \item $15.3\times$ gas savings for $N{=}8$ enterprises, scaling to
    $30.5\times$ at $N{=}16$.
  \item Sub-15\,s aggregation time with constant memory.
  \item Requires audit of the Sonobe library and maturation of
    ProtoGalaxy implementations.
  \item Per-circuit Groth16 trusted setup for the decider circuit; can
    be amortized since the decider circuit is fixed.
\end{itemize}

\subsection{Anti-Confirmation Bias Assessment}

We identify the following conditions under which the primary
recommendation (folding + Groth16 decider) would be weakened:

\begin{enumerate}
  \item If the ProtoGalaxy verifier circuit exceeds 100K constraints
    (actual: ${\sim}11$K), accumulation-based approaches would be
    comparably efficient.
  \item If the Groth16 decider trusted setup proves operationally
    burdensome (per-circuit ceremony), the universal-setup advantage of
    binary tree accumulation becomes more valuable.
  \item If the Sonobe library reveals security vulnerabilities in audit,
    the fallback to binary tree accumulation is immediate and
    well-understood.
  \item For $N \leq 3$ enterprises, the engineering complexity of any
    aggregation strategy may not justify the gas savings, particularly
    on a zero-fee network where gas cost is computational rather than
    monetary.
\end{enumerate}

\textbf{Steelman for no aggregation:} On the Basis Network L1 with zero
gas fees, the monetary motivation for aggregation is absent. The
throughput argument (block capacity) only becomes binding at
$N > 10$--$15$ enterprises. For the initial deployment with fewer than
5 enterprises, individual verification is simpler, equally correct, and
imposes no aggregation latency.
